% Full OECP Paper in LaTeX
\documentclass[12pt]{article}
\usepackage{amsmath, amssymb, graphicx, hyperref, setspace}
\onehalfspacing

\title{Open Economy Carbon Pricing: A Trade-Focused Framework for Efficient, Politically Feasible, and Rapid Decarbonisation}
\author{Stephen J. Stretton}
\date{Full Draft}

\begin{document}
\maketitle

% ==================== ABSTRACT ====================
\begin{abstract}
Carbon pricing remains too weak in most countries to deliver rapid decarbonisation. We argue this weakness is structural, arising from three constraints: (1) the trade-exposure constraint, (2) the political-economy constraint, and (3) the primary-to-quaternary structure of energy transition. Conventional carbon pricing (CCP) applies a source-based tax to domestic production, creating competitiveness risks and inviting a race to the bottom similar to corporate income taxation. Open Economy Carbon Pricing (OECP) restructures carbon pricing into a destination-based framework, pairing a production-side feebate with a consumption-based carbon charge that applies symmetrically to imports and domestic goods while zero-rating exports. This enables a race to the top akin to the rise of VAT. We present the simplest integrated implementation mechanism using CBAM logic and common external tariff structures. Appendices provide a simple trade model and a simple political-feasibility model.
\end{abstract}


\tableofcontents
% ==================== INTRODUCTION ====================
\section{Introduction}
Carbon pricing has not achieved the levels or pace of decarbonisation consistent with climate goals. This is not due to inadequate political ambition or weak modelling but to three underlying structural constraints.

First, CCP is a \textbf{source-based tax}. It applies to emissions at the point of production. Like corporate income taxes, source-based carbon taxes create competitiveness pressures: raising the tax increases domestic costs and encourages production to shift abroad. The result is a \textbf{race to the bottom}.

Second, CCP faces \textbf{political-economy limits}. Costs are concentrated on producers and salient to households. Voter and industry resistance sharply caps feasible carbon prices.

Third, CCP is structurally \textbf{too weak to accelerate primary-sector transformation}. Steel, cement, fuels, and electricity require capital turnover and technology switching, not marginal adjustments. A pure carbon tax would need to be very high to trigger switching, and such a tax is politically impossible.

Open Economy Carbon Pricing (OECP) resolves these constraints by separating the incentive and revenue functions of carbon pricing. OECP consists of:

\begin{enumerate}
    \item A \textbf{revenue-neutral production feebate} that changes relative costs between high- and low-emitting producers without raising average cost.
    \item A \textbf{destination-based consumption charge} applied to all goods consumed domestically, regardless of origin, with exports zero-rated.
\end{enumerate}

This structure mirrors VAT, which allowed countries to raise tax rates without damaging competitiveness. OECP creates the same dynamic for carbon pricing. This paper focuses primarily on the \textbf{trade argument} and presents a simple implementation suitable for rapid adoption.

% ==================== THREE ARGUMENTS ====================
\section{Three Structural Arguments for OECP}

\subsection{Trade Argument (Primary Focus)}
CCP taxes emissions where production occurs. This raises producers' costs directly. When one country raises its carbon price, its exporters lose competitiveness relative to producers abroad. Imports also gain market share. Policymakers perceive this immediately as a threat to jobs, tax bases, and industrial capability.

Thus CCP generates \textbf{international carbon-tax competition}. Each country fears raising its carbon price because production will shift abroad. This strategic interaction mirrors the deterioration seen in corporate income tax rates worldwide.

OECP addresses this by shifting carbon pricing to a \textbf{destination basis}. The consumption charge applies equally to imported and domestic goods, while exports are zero-rated. Producers face incentives to decarbonise through the feebate but do not face an absolute increase in costs that would harm competitiveness.

This reproduces the same "race to the top" dynamic seen with VAT. Because VAT exempts exports and taxes imports, governments can raise VAT rates without losing competitiveness. OECP applies this logic to carbon.

\subsection{Political-Economy Argument}
The political problem of carbon pricing is simple: most voters actively dislike higher energy prices, and most firms dislike higher production costs. CCP delivers both.

Under OECP, these two constraints are relaxed:

\begin{enumerate}
    \item The \textbf{production feebate} recycles revenue within each sector, eliminating most net costs for producers.
    \item The \textbf{consumption charge} generates revenue that can be returned as household dividends, lowering political resistance.
\end{enumerate}

Appendix B presents a simple model showing that the feasible carbon price under CCP is a fraction of that under OECP.

\subsection{Accelerated-Transition Argument (Primary to Quaternary Structure)}
Decarbonisation begins with the primary sectors (steel, cement, fuels, electricity). These sectors are capital-intensive, with long-lived assets and limited short-run substitution. A carbon tax must be very high to shift them.

A feebate, by contrast, acts through \textbf{relative incentives}. It rewards clean production and penalises dirty production regardless of average cost, accelerating switching once the differential meets the technology cost gap.

CCP cannot achieve this without raising export prices. OECP achieves this without harming competitiveness.

% ==================== IMPLEMENTATION ====================
\section{Implementation: A Unified and Simple OECP Mechanism}
The OECP framework can be implemented using the simplest possible tools. This section outlines the core instruments.

\subsection{The Production Feebate}
Let $t$ be the carbon incentive rate. Let $e_i$ and $\bar e$ denote individual and benchmark emissions intensities. The feebate is:

\begin{equation}
F_i = t (e_i - \bar e).
\end{equation}

This ensures strong incentives without raising average prices. Export competitiveness is preserved.

\subsection{The Consumption Charge}
The consumption charge applies to all goods consumed domestically:

\begin{equation}
C_j = t m_j,
\end{equation}
where $m_j$ is embodied emissions.

Imports are taxed at the border; domestic goods at the point of sale. Exports are zero-rated, as under VAT.

\subsection{Import Adjustments: CBAM Logic Without Complexity}
Imports face a charge equal to $t m_j$. Embodied emissions may be:
\begin{itemize}
    \item Verified by producers; or
    \item Determined using regulator-issued default values.
\end{itemize}

Unlike CBAM, OECP does not require tracking foreign carbon prices or foreign free allocations.

\subsection{Export Adjustments: VAT Logic}
Exports are zero-rated. Producers retain feebate incentives but face no consumption charge. This preserves export competitiveness.

\subsection{Common External Tariff Option}
A simple extension adds a base tariff plus a carbon adjustment. Imports face:

\begin{enumerate}
    \item A base tariff, and
    \item A carbon surcharge equal to $t m_j$.
\end{enumerate}

Low-carbon exporters may receive reduced surcharges after verification.

% ==================== APPENDICES ====================
\appendix

\section{Appendix A: Simple Trade Model}
Consider a small open economy with two technologies: dirty (D) and clean (C).

Dirty has cost $c_D$ and emissions $e_D$. Clean has cost $c_C = c_D + \Delta c$ and emissions $e_C$.

World price is $P$. Under CCP, exporters face a domestic price:
\begin{equation}
P_D = P + t e_D.
\end{equation}

Under OECP, export prices remain $P$.

The switching condition under CCP is:
\begin{equation}
P - c_D - t e_D \le P - c_C - t e_C,
\end{equation}
which includes competitiveness losses.

Under OECP, switching is driven only by the feebate:
\begin{equation}
c_D - c_C = t (e_D - e_C).
\end{equation}

Thus OECP restores first-best incentives without raising export prices.

\section*{Appendix B: Full Political Feasibility Model}


This appendix develops a minimal political-economy model illustrating why the feasible carbon price under Conventional Carbon Pricing (CCP) is structurally lower than under Open Economy Carbon Pricing (OECP). The purpose is not to model political dynamics in full but to provide a simple representation of the binding constraints.


%%%%%%%%%%%%%%%%%%%%%%%%%%%%%%%%%%%%%%%%%%%%%%%%%%%%%%%%%%%%%%
\subsection*{B.1 Setup}


There are two politically relevant groups:
\begin{itemize}
\item Firms (especially trade-exposed manufacturers)
\item Voters (households)
\end{itemize}


Each group experiences a perceived burden from carbon pricing. Perceived burdens may differ sharply from physical burdens because political salience, media narratives, and behavioural biases shape reactions.


Let:
\begin{itemize}
\item $t$ = effective carbon price level
\item $a_F$ = burden intensity perceived by firms per unit of carbon price
\item $a_V$ = burden intensity perceived by voters per unit of carbon price
\item $\rho_F$ = share of firm burden that is rebated/neutralised
\item $\rho_V$ = share of voter burden that is rebated/neutralised
\item $K_F$ = tolerance threshold of firms
\item $K_V$ = tolerance threshold of voters
\end{itemize}


The political system is stable only if both constraints are satisfied.


%%%%%%%%%%%%%%%%%%%%%%%%%%%%%%%%%%%%%%%%%%%%%%%%%%%%%%%%%%%%%%
\subsection*{B.2 Political Acceptance Conditions}


\paragraph{Firm acceptance.}
Firms accept the carbon price if the net perceived burden is below their tolerance:
\begin{equation}
a_F t (1 - \rho_F) \le K_F.
\end{equation}


\paragraph{Voter acceptance.}
Voters accept the carbon price if their net perceived burden is below their tolerance:
\begin{equation}
a_V t (1 - \rho_V) \le K_V.
\end{equation}


Thus the maximum feasible carbon price is:
\begin{equation}
t^{max} = \min \left( \frac{K_F}{a_F (1 - \rho_F)}, \; \frac{K_V}{a_V (1 - \rho_V)} \right).
\end{equation}


Under CCP, both rebate rates are close to zero. Under OECP, both rebate rates rise significantly.


%%%%%%%%%%%%%%%%%%%%%%%%%%%%%%%%%%%%%%%%%%%%%%%%%%%%%%%%%%%%%%
\subsection*{B.3 CCP Case: Binding Low Ceiling}


Under CCP (production carbon tax):
\begin{itemize}
\item Firms face the full cost of the carbon tax: $\rho_F^{CCP} \approx 0$.
\item Voters face higher energy prices with limited compensation: $\rho_V^{CCP} \ll 1$.
\end{itemize}
\end{document}